%%%%%%%%%%%%%%%%%
% CV tailored for Intheair / R&D GPR, Geophysics & Drone Data Fusion
%%%%%%%%%%%%%%%%%

\documentclass[8pt,a4paper,ragged2e,withhyper]{../_shared/altacv}

\geometry{left=1.25cm,right=1.25cm,top=.5cm,bottom=1.cm,columnsep=1.2cm}

\usepackage{paracol}
\usepackage{fontawesome5}
\usepackage{hyperref}

\iftutex 
  \setmainfont{Roboto Slab}
  \setsansfont{Lato}
  \renewcommand{\familydefault}{\sfdefault}
\else
  \usepackage[rm]{roboto}
  \usepackage[defaultsans]{lato}
  \renewcommand{\familydefault}{\sfdefault}
\fi

\definecolor{SlateGrey}{HTML}{2E2E2E}
\definecolor{LightGrey}{HTML}{666666}
\definecolor{DarkPastelRed}{HTML}{8AC0D9}
\definecolor{PastelRed}{HTML}{00B0DA}
\definecolor{GoldenEarth}{HTML}{B4AD0C}

\colorlet{name}{black}
\colorlet{tagline}{PastelRed}
\colorlet{heading}{DarkPastelRed}
\colorlet{headingrule}{GoldenEarth}
\colorlet{subheading}{PastelRed}
\colorlet{accent}{PastelRed}
\colorlet{emphasis}{SlateGrey}
\colorlet{body}{LightGrey}

\definecolor{violet}{HTML}{5A2582}
\definecolor{rouge}{HTML}{F53B3B}
\definecolor{noir}{HTML}{00232C}
\definecolor{bleu}{HTML}{00B0DA}
\definecolor{rose}{HTML}{F831A0}
\definecolor{jaune}{HTML}{FFEC00}
\definecolor{vert}{HTML}{34AD6C}

\renewcommand{\namefont}{\Huge\rmfamily\bfseries}
\renewcommand{\personalinfofont}{\footnotesize}
\renewcommand{\cvsectionfont}{\LARGE\rmfamily\bfseries}
\renewcommand{\cvsubsectionfont}{\large\bfseries}
\renewcommand{\cvItemMarker}{{\small\textbullet}}
\renewcommand{\cvRatingMarker}{\faCircle}

\input{../_shared/pubs-num.tex}
\addbibresource{../_shared/sample.bib}

\begin{document}

\name{BOILEAU Guillaume}
\tagline{Ingénieur R\&D Physique, Signal \& IA | Détection faible SNR, Python, Fusion de données, Validation terrain}

\photoL{2.8cm}{../_shared/moi}

\personalinfo{%
  \email{guillaume.boileaupro@gmail.com}
  \phone{+33 6 59 94 85 84}
  \location{Alpes-Maritimes, FRANCE}
  \linkedin{guillaume-boileau-phd-891b81265}
  \researchgate{Guillaume-Boileau-2}
  \orcid{0000-0002-3576-6968}
}

\makecvheader
\vspace{-0.3cm}

\AtBeginEnvironment{itemize}{\small}
\columnratio{0.64}

\begin{paracol}{2}

\cvsection{Experience}

\cvevent{\textbf{Software Engineer -- System Integration, Data Flows \& Operational Validation}}{VEV Platform Services France}{2025 -- 2026}{Valbonne, France}

\textbf{Context:} Industrial software environment for an EV charging CPMS platform connected to operational charging stations, field equipment and hardware/software interfaces.

\textbf{Objective:} Support reliable operation of connected infrastructure through software integration, data-flow checks, operational monitoring, anomaly analysis and validation of technical behaviours.

\begin{itemize}
  \item \textbf{Operational data flows:} Implemented and monitored FTP-based exchanges between software services and external systems.
  \item \textbf{Anomaly analysis:} Investigated interrupted exchanges, inconsistent flows and defects affecting service behaviour.
  \item \textbf{Validation:} Performed checks in production-like environments to assess continuity, expected behaviour and data consistency.
  \item \textbf{Software support:} Maintained TypeScript, Angular and Node.js components and documented technical findings.
\end{itemize}

\divider

\cvevent{\textbf{Senior R\&D Engineer -- Signal Processing, Modelling \& Validation Workflows}}{CNES, Laboratoire Lagrange, Observatoire de la Côte d'Azur}{2023 -- 2025}{Nice, France}

\textbf{Context:} R\&D activities for the LISA ESA/NASA space mission, involving physical modelling, weak-signal analysis, simulation, software chains, data-processing workflows and validation campaigns.

\textbf{Objective:} Develop and validate scientific analysis workflows able to compare complex models, process large datasets and assess weak physical signatures under strong noise and consistency constraints.

\begin{itemize}
  \item \textbf{Weak-signal analysis:} Worked on low signal-to-noise physical backgrounds and complex superposed signals.
  \item \textbf{Signal processing:} Developed methods for spectral interpretation, model comparison and performance assessment.
  \item \textbf{Python scientific platform:} Developed CATpyLISA for large-scale model integration, comparison and validation.
  \textcolor{DarkPastelRed}{\href{https://gitlab.oca.eu/gboileau/lisa_l3_tools}{\bf GitLab: CATpyLISA}}
  \item \textbf{Validation campaigns:} Prepared and analysed comparison campaigns on simulated outputs and heterogeneous models.
  \item \textbf{Reproducible workflows:} Used Python, Linux, Git/GitLab and Jupyter-style environments for traceable analyses.
  \item \textbf{Collaboration:} Worked with international contributors across software, systems, modelling and instrumentation.
\end{itemize}

\divider

\cvevent{\textbf{Data Scientist -- Noise Analysis, Statistical Modelling \& Performance Validation}}{Universiteit Antwerpen, Belgium}{2021 -- 2023}{Antwerp, Belgium}

\textbf{Context:} R\&D work for Einstein Telescope and next-generation gravitational-wave detector studies, combining environmental noise analysis, detector-configuration studies, statistical modelling and weak-signal detection.

\textbf{Objective:} Analyse the impact of environmental noise on detector sensitivity, develop statistical workflows for different detector configurations and support technical recommendations on noise budgets, geometry and experimental validation.

\begin{itemize}
  \item \textbf{Geophysical data analysis:} Used seismometer measurements to analyse propagation effects and spatial correlations.
  \item \textbf{Noise studies:} Investigated seismic and Newtonian noise impacts on detector sensitivity.
  \item \textbf{Experimental preparation:} Contributed to protocols, sensor selection and deployment preparation for low-noise campaigns.
  \item \textbf{Scientific Python workflows:} Developed reproducible pipelines for signal analysis, statistical modelling and validation.
\end{itemize}

\divider

\cvevent{\textbf{Junior R\&D Engineer -- Simulation, Signal Diagnostics \& Scientific Computing}}{ARTEMIS, Observatoire de la Côte d'Azur}{2018 -- 2021}{Nice, France}

\textbf{Context:} Software and analysis activities for gravitational-wave mission preparation, involving simulation, signal decomposition, diagnostic analysis and validation of complex physical environments.

\textbf{Objective:} Develop robust analysis workflows for weak-signal detection, simulation, uncertainty studies and interpretation of complex datasets.

\begin{itemize}
  \item \textbf{Signal decomposition:} Developed methods to separate overlapping low-SNR signals.
  \item \textbf{Simulation:} Built approaches for complex signal environments and uncertainty studies.
  \item \textbf{Diagnostics:} Investigated noise sources through cross-sensor correlation analyses.
  \item \textbf{Scientific Python:} Developed Python-based analysis, modelling and reproducible workflows.
\end{itemize}

\switchcolumn

\cvsection{Profile for Intheair}

\textbf{R\&D engineer and PhD physicist with experience in signal processing, noise analysis, scientific Python and validation of complex physical models.} Background developed across space systems, precision instrumentation, connected infrastructure and scientific software.

Transferable strengths for this role: physical modelling, low-SNR analysis, reproducible data workflows, experimental reasoning and validation of complex measurements.

\begin{itemize}
  \item Strong background in physics, modelling and interpretation of complex signals.
  \item Experience in low-SNR analysis, noise diagnostics and performance assessment.
  \item Python development for reproducible data-processing and validation workflows.
  \item Used to working across data, software, physical systems and field constraints.
  \item Able to ramp up quickly on GPR, geophysics and drone-based data workflows.
\end{itemize}

\cvsection{Languages}

\cvskill{English}{4}
\divider
\cvskill{Spanish}{2}
\divider

\medskip

\cvsection{Education}

\cvevent{PhD in Physics}{Université Côte d'Azur}{2018 -- 2021}{Nice, France}

\divider

\cvevent{Selective Graduate Program in Fundamental Physics (Magistère)}{Université Paris-Saclay}{2016 -- 2018}{Orsay, France}

\divider

\cvevent{MSc in Astronomy and Astrophysics}{Observatoire de Paris / Université Paris-Saclay}{2017 -- 2018}{Paris, France}

\divider

\cvevent{BSc in Fundamental Physics}{Université Paris-Saclay}{2015 -- 2016}{Orsay, France}

\cvsection{Professional Training}

\cvevent{Machine Learning in Python with scikit-learn}{FUN MOOC -- Inria}{2026}{France Université Numérique}

\small{
Predictive modelling, supervised learning, model evaluation, cross-validation, metrics, pipelines and practical machine learning workflows with Python and scikit-learn.
}

\divider

\cvevent{Supervised Machine Learning (Regression \& Classification)}{DeepLearning.AI -- Andrew Ng}{2020}{Coursera}

\divider

\cvevent{Recherche reproductible : principes méthodologiques pour une science transparente}{FUN MOOC -- Inria}{2025}{France Université Numérique}

\divider

\cvevent{Continuous Professional Training -- Software, Data, AI and Systems}{OpenClassrooms}{2024 -- 2026}{}

\small{
Python, SQL, Machine Learning, AI, Linux, JavaScript, TypeScript, Angular, Node.js, Django, REST APIs, C/C++, web development, algorithms and software engineering fundamentals.
}

\end{paracol}

\cvsection{Projects}

\cvevent{LISA / CATpyLISA -- Weak-Signal Analysis, Model Validation \& Scientific Python}{ESA/NASA Space Mission -- Large-scale International Collaboration}{2023 -- 2025}{Nice, France}

Python workflows for weak-signal analysis, model comparison and validation of complex physical data.

\divider

\cvevent{Noise Diagnostics \& Performance Analysis for Precision Instrumentation}{Signal Processing / Statistical Modelling / Environmental Effects}{2021 -- 2023}{Antwerp, Belgium}

Statistical and signal-analysis workflows to identify environmental and instrumental noise sources.

\divider

\cvevent{VEV -- Connected Infrastructure, Operational Data Flows \& System Validation}{Industrial Software / Field Equipment Interfaces}{2025 -- 2026}{Valbonne, France}

Software integration, operational monitoring and validation for a platform connected to field infrastructure.

\cvsection{Strengths}

\vspace{-.3cm}

\cvsubsection{Signal processing \& weak-signal detection}

{\small
\cvtag{Signal Processing}
\cvtag{Weak-Signal Detection}
\cvtag{Low SNR Analysis}
\cvtag{Noise Reduction}
\cvtag{Spectral Analysis}
\cvtag{Filtering}
\cvtag{Correlation Analysis}
\cvtag{Signal Decomposition}
\cvtag{Anomaly Detection}
\cvtag{Performance Analysis}
\cvtag{Uncertainty Analysis}
\cvtag{Model Comparison}
}

\divider\smallskip
\vspace{-.3cm}

\cvsubsection{Physics, radar \& geophysics ramp-up}

{\small
\cvtag{Physics}
\cvtag{Electromagnetic Modelling Awareness}
\cvtag{Propagation Effects}
\cvtag{Radar / GPR Ramp-Up}
\cvtag{Radargram Interpretation Ramp-Up}
\cvtag{Near-Surface Geophysics Ramp-Up}
\cvtag{Subsurface Mapping Ramp-Up}
\cvtag{Sensor Calibration Logic}
\cvtag{Experimental Protocols}
\cvtag{Field Validation}
\cvtag{Physical Modelling}
}

\divider\smallskip
\vspace{-.3cm}

\cvsubsection{Python, data science \& AI}

{\small
\cvtag{Python}
\cvtag{NumPy}
\cvtag{SciPy}
\cvtag{pandas}
\cvtag{scikit-learn}
\cvtag{Matplotlib}
\cvtag{Jupyter}
\cvtag{Machine Learning}
\cvtag{Data Cleaning}
\cvtag{Feature Extraction}
\cvtag{Classification}
\cvtag{Regression}
\cvtag{Scientific Computing}
\cvtag{Reproducible Workflows}
}

\divider\smallskip
\vspace{-.3cm}

\cvsubsection{Geospatial \& multi-sensor data}

{\small
\cvtag{Data Fusion}
\cvtag{Multi-Source Data}
\cvtag{Geospatial Data Ramp-Up}
\cvtag{GIS Ramp-Up}
\cvtag{QGIS Ramp-Up}
\cvtag{LiDAR Awareness}
\cvtag{Photogrammetry Awareness}
\cvtag{GNSS RTK Awareness}
\cvtag{Digital Twins Awareness}
\cvtag{Mapping Workflows}
\cvtag{Coordinate Systems Ramp-Up}
\cvtag{Data Georeferencing Ramp-Up}
}

\divider\smallskip
\vspace{-.3cm}

\cvsubsection{Validation, software \& industrialisation}

{\small
\cvtag{R\&D Validation}
\cvtag{Validation Strategy}
\cvtag{Test Protocols}
\cvtag{Field-Test Logic}
\cvtag{Technical Reporting}
\cvtag{Traceability}
\cvtag{Quality Follow-Up}
\cvtag{Linux}
\cvtag{Git / GitLab}
\cvtag{GitHub}
\cvtag{Docker}
\cvtag{CI/CD}
\cvtag{Bash}
\cvtag{SQL}
\cvtag{C}
\cvtag{C++}
}

\divider\smallskip
\vspace{-.3cm}

\cvsubsection{Collaboration \& R\&D delivery}

{\small
\cvtag{Technical English}
\cvtag{International Environment}
\cvtag{Stakeholder Coordination}
\cvtag{Cross-Functional Collaboration}
\cvtag{Agile}
\cvtag{Kanban}
\cvtag{Backlog Follow-Up}
\cvtag{Jira}
\cvtag{Confluence}
\cvtag{Zenhub}
\cvtag{Research-to-Product}
\cvtag{Deeptech}
\cvtag{Industrial Software}
\cvtag{Field Constraints}
\cvtag{Autonomy}
}

\end{document}
